% -*- coding: utf-8 -*-
% !TEX program = xelatex
\documentclass[plain,noanswer,medmath]{../../template/jnuexam}
\usepackage{../../template/kysx}

\begin{document}


\examtitle{1996}{4}


\makepart{填空题}{本题共5小题，每小题3分，满分15分}


\begin{problem}
设方程$x = y^y$确定$y$是$x$的函数，则$\dy =$ \fillin{}．
\end{problem}


\begin{problem}
设$\int x f(x)\dx = \arcsin x + C$，则$\int \frac{1}{f(x)} \dx =$ \fillin{}．
\end{problem}


\begin{problem}
设$\left(x_0,y_0 \right)$是抛物线$y = ax^2 + bx + c$上的一点，若在该点的切线过原点，则系数应满足的关系是 \fillin{}．
\end{problem}


\begin{problem}
设
$$A = \begin{pmatrix}
            1 &           1 &           1 & \cdots & 1 \\
          a_1 &         a_2 &         a_3 & \cdots & a_n \\
        a_1^2 &       a_2^2 &       a_3^2 & \cdots & a_n^2 \\
       \vdots &      \vdots &      \vdots &        & \vdots \\
  a_1^{n - 1} & a_2^{n - 1} & a_3^{n - 1} & \cdots & a_n^{n - 1}
\end{pmatrix},\quad X = \begin{pmatrix}
  x_1 \\
  x_2 \\
  x_3 \\
  \vdots \\
  x_n
\end{pmatrix},\quad B = \begin{pmatrix}
  1 \\
  1 \\
  1 \\
  \vdots \\
  1
\end{pmatrix},$$
其中$a_i \ne a_j$（$i \ne j;i,j = 1,2, \cdots ,n$）．则线性方程组$A^T X = B$的解是 \fillin{}．
\end{problem}


\begin{problem}
设由来自正态总体$X\sim N(\mu ,0.9^2)$容量为9的简单随机样本，得样本均值$\overline{X}= 5$，
则未知参数$\mu$的置信度为0.95的置信区间为 \fillin{}．
\end{problem}


\makepart{选择题}{本题共5小题，每小题3分，满分15分}


\begin{problem}
累次积分$\int_0^{\frac{\pi}{2}} \d\theta \int_0^{\cos \theta} f(r\cos \theta ,r\sin \theta) r\dr$
可以写成\pickout{}
\begin{abcd}
\item $\int_0^1\dy\int_0^{\sqrt{y - y^2}} f(x,y) \dx$．				
\item $\int_0^1\dy\int_0^{\sqrt{1 - y^2}} f(x,y) \dx$．
\item $\int_0^1\dx\int_0^1f(x,y) \dy$．					
\item $\int_0^1\dx\int_0^{\sqrt{x - x^2}} f(x,y) \dy$．
\end{abcd}
\end{problem}


\begin{problem}
下述各选项正确的是\pickout{}
\begin{abcd}
\item 若$\sum_{n = 1}^{\infty}u_n^2$和$\sum_{n = 1}^{\infty}v_n^2$都收敛，
      则$\sum_{n = 1}^{\infty}(u_n^{} + v_n)^2$收敛．
\item $\sum_{n = 1}^{\infty}\left| u_n^{}v_n \right|$收敛，
      则$\sum_{n = 1}^{\infty}u_n^2$与$\sum_{n = 1}^{\infty}v_n^2$都收敛．
\item 若正项级数$\sum_{n = 1}^{\infty}u_n^{}$发散，则$u_n^{} \ge \frac{1}{n}$．
\item 若级数$\sum_{n = 1}^{\infty}u_n^{}$收敛，且$u_n^{} \ge v_n$（$n = 1,2, \cdots$），
      则级数$\sum_{n = 1}^{\infty}v_n^{}$也收敛．
\end{abcd}
\end{problem}


\begin{problem}
设$n$阶矩阵$A$非奇异（$n \ge 2$），$A^*$是矩阵$A$的伴随矩阵，则\pickout{}
\begin{abcd}
\item $(A^*)^* = |A|^{n - 1} A$．		
\item $(A^*)^* = |A|^{n + 1} A$．
\item $(A^*)^* = |A|^{n - 2} A$．		
\item $(A^*)^* = |A|^{n + 2} A$．
\end{abcd}
\end{problem}


\begin{problem}
设有任意两个$n$维向量组$\alpha_1, \cdots ,\alpha_m$和$\beta_1, \cdots ,\beta_m$,
若存在两组不全为零的数$\lambda_1, \cdots ,\lambda_m$和$k_1, \cdots ,k_m$，使
$$(\lambda_1 + k_1)\alpha_1 + \cdots + (\lambda_m + k_m)\alpha_m + (\lambda_1 - k_1)\beta_1
   + \cdots + (\lambda_m - k_m)\beta_m = 0,$$
则\pickout{}
\begin{abcd}
\item $\alpha_1, \cdots ,\alpha_m$和$\beta_1, \cdots ,\beta_m$都线性相关．
\item $\alpha_1, \cdots ,\alpha_m$和$\beta_1, \cdots ,\beta_m$都线性无关．
\item $\alpha_1+\beta_1, \cdots, \alpha_m+\beta_m, \alpha_1-\beta_1, \cdots, \alpha_m-\beta_m$线性无关．
\item $\alpha_1+\beta_1, \cdots, \alpha_m+\beta_m, \alpha_1-\beta_1, \cdots, \alpha_m-\beta_m$线性相关．
\end{abcd}
\end{problem}


\begin{problem}
已知$0 < P(B) < 1$且$P[(A_1 + A_2)|B] = P(A_1|B) + P(A_2|B)$，则下列选项成立的是\pickout{}
\begin{abcd}
\item $P[(A_1 + A_2)|\overline{B}] = P(A_1|\overline{B}) + P(A_2|\overline{B})$．
\item $P(A_1 B + A_2 B) = P(A_1 B) + P(A_2 B)$．
\item $P(A_1 + A_2) = P(A_1|B) + P(A_2|B)$．
\item $P(B) = P(A_1)P(B|A_1) + P(A_2)P(B|A_2)$．
\end{abcd}
\end{problem}


\makepart{}{本题满分6分}

\begin{problem}
设$f(x) = \begin{cases}
  \frac{g(x) - \e^{- x}}{x}, & x \ne 0, \\
                          0, & x = 0,
\end{cases}$其中$g(x)$有二阶连续导数，且$g(0) = 1,g'(0) = - 1$．\par
\begin{enumerate*}
\item 求$f'(x)$；
\item 讨论$f'(x)$在$(-\infty , +\infty)$上的连续性．
\end{enumerate*}
\end{problem}


\makepart{}{本题满分6分}

\begin{problem}
设函数$z = f(u)$，方程$u = \varphi(u) + \int_y^x p(t)\dt$确定$u$是$x,y$的函数，
其中$f(u),\varphi(u)$可微；$p(t)$,$\varphi'(u)$连续，且$\varphi'(u) \ne 1$．
求$p(y)\frac{\pdz}{\pdx} + p(x)\frac{\pdz}{\pdy}$．
\end{problem}


\makepart{}{本题满分6分}

\begin{problem}
计算$\int_0^{+\infty} \frac{x\e^{- x}}{(1 + \e^{- x})^2} \dx$．
\end{problem}


\makepart{}{本题满分5分}

\begin{problem}
设$f(x)$在区间$[0,1]$上可微，且满足条件$f(1) = 2\int_0^{\frac{1}{2}} xf(x)\dx$．
试证：存在$\xi \in(0,1)$使$f(\xi) + \xi f'(\xi) = 0$．
\end{problem}


\makepart{}{本题满分6分}

\begin{problem}
设某种商品的单价为$p$时，售出的商品数量$Q$可以表示成$Q = \frac{a}{p + b} - c$,
其中$a$, $b$, $c$均为正数，且$a > bc$．
\begin{enumerate}
\item 求$p$在何范围变化时，使相应销售额增加或减少．
\item 要使销售额最大，商品单价$p$应取何值？最大销售额是多少？
\end{enumerate}
\end{problem}


\makepart{}{本题满分6分}

\begin{problem}
求微分方程$\frac{\dy}{\dx} = \frac{y - \sqrt{x^2 + y^2}}{x}$的通解．
\end{problem}


\makepart{}{本题满分8分}

\begin{problem}
设矩阵$A = \begin{pmatrix}
  0 & 1 & 0 & 0 \\
  1 & 0 & 0 & 0 \\
  0 & 0 & y & 1 \\
  0 & 0 & 1 & 2
\end{pmatrix}$．
\begin{enumerate}
\item 已知$A$的一个特征值为$3$，试求$y$；
\item 求矩阵$P$，使$(AP)^T(AP)$为对角矩阵．
\end{enumerate}
\end{problem}


\makepart{}{本题满分8分}

\begin{problem}
设向量$\alpha_1,\alpha_2, \cdots ,\alpha_t$是齐次线性方程组$AX = 0$的一个基础解系，
向量$\beta$不是方程组$AX = 0$的解，即$A\beta \ne 0$．
试证明：向量组$\beta ,\beta + \alpha_1,\beta + \alpha_2, \cdots ,\beta + \alpha_t$线性无关．
\end{problem}


\makepart{}{本题满分7分}

\begin{problem}
假设一部机器在一天内发生故障的概率为$0.2$，机器发生故障时全天停止工作，
若一周$5$个工作日里无故障，可获利润$10$万元；
发生一次故障仍可获得利润$5$万元；发生两次故障所获利润$0$元；发生三次或三次以上故障就要亏损$2$万元．
求一周内期望利润是多少?
\end{problem}


\makepart{}{本题满分6分}

\begin{problem}
考虑一元二次方程$x^2 + Bx + C = 0$，其中$B,C$分别是将一枚色子(骰子)接连掷两次先后出现的点数．
求该方程有实根的概率$p$和有重根的概率$q$．
\end{problem}


\makepart{}{本题满分6分}

\begin{problem}
假设$X_1,X_2, \cdots ,X_n$是来自总体X的简单随机样本；已知$EX^k = a_k(k = 1,2,3,4)$．
证明：当$n$充分大时，随机变量$Z_n = \frac{1}{n}\sum_{i = 1}^n X_i^2$近似服从正态分布，并指出其分布参数．
\end{problem}


\end{document}
